% !TeX root = ../main.tex
\begin{figure}[H]
\centering
\small
\begin{mathpar}
 \inferrule*[right=Refl]{ }{\tau<:\tau}
 \and
 \inferrule*[right=Trans]{\tau<:\sigma\\\sigma<:\rho}{\tau<:\rho}
 \and
 \inferrule*[right=Real-Sub]{ }{\tyR\G<:\tyR\E}
 \and
 \inferrule*[right=Arrow]{\sigma_1<:\tau_1\\\tau_2<:\sigma_2}
   {\tau_1\to\tau_2<:\sigma_1\to\sigma_2}
 \and
 \inferrule*[right=Product]{\tau_1<:\sigma_1\\\tau_2<:\sigma_2}
   {\tau_1\times\tau_2<:\sigma_1\times\sigma_2}
 \and
 \inferrule*[right=Sum]{\tau_1<:\sigma_1\\\tau_2<:\sigma_2}
   {\tau_1+\tau_2<:\sigma_1+\sigma_2}
 \and
 \inferrule*[right=List]{\tau<:\sigma}{\tyL\tau<:\tyL\sigma}
\end{mathpar}
\begin{mathpar}
 \inferrule*[right=Sub]{\typing e\tau\\\tau<:\sigma}{\typing e\sigma}
 \and
 \inferrule*[right=Real]{r\in\Real}{\typing r{\tyR m}}
 \and
 \inferrule*[right=Reject]{ }{\typing{\kw{reject}}\tau}
 \and
 \inferrule*[right=Neg]{\typing e{\tyR m}}{\typing{-e}{\tyR m}}
 \and
 \inferrule*[right=Add]{\typing{e_1}{\tyR m}\\\typing{e_2}{\tyR m}}
   {\typing{e_1+e_2}{\tyR m}}
 \and
 \inferrule*[right=Mul]{\typing{e_1}{\tyR\G}\\\typing{e_2}{\tyR m}}
   {\typing{e_1\cdot e_2}{\tyR m}}
 \and
 \inferrule*[right=Div]{\typing{e_1}{\tyR m}\\\typing{e_2}{\tyR\G}}
   {\typing{e_1/e_2}{\tyR m}}
 \and
 \inferrule*[right=Cmp]{\typing{e_1}{\tyR\G}\\\typing{e_2}{\tyR\G}}
   {\typing{e_1<e_2}{\tyB}}
\end{mathpar}
\begin{mathpar}
 \inferrule*[right=Uniform]{\typing a{\tyR m}\\\typing b{\tyR m}}
   {\typing{\draw{uniform}\alpha{a,b}}{\tyR m}}
 \and
 \inferrule*[right=Gaussian]{\typing u{\tyR m}\\\typing v{\tyR\G}}
   {\typing{\draw{gaussian}\alpha{u,v}}{\tyR m}}
 \and
 \inferrule*[right=Poisson]{\typing\lambda{\tyR m}}
   {\typing{\draw{poisson}\alpha\lambda}{\tyR m}}
 \and
 \inferrule*[right=Exponential]{\typing\lambda{\tyR\G}}
   {\typing{\draw{exponential}\alpha\lambda}{\tyR m}}
 \and
 \inferrule*[right=Bernoulli]{\typing p{\tyR m}}
   {\typing{\draw{bernoulli}\alpha p}{\tyR m}}
 \and
 \inferrule*[right=Discrete]{\typing p{\tyL{\tyR m}}}
   {\typing{\draw{discrete}\alpha p}{\tyR m}}
 \and
 \inferrule*[right=Beta]{\typing a{\tyR\G}\\\typing b{\tyR\G}}
   {\typing{\draw{beta}\alpha{a,b}}{\tyR m}}
 \and
 \inferrule*[right=Gamma]{\typing k{\tyR m}\\\typing\lambda{\tyR\G}}
   {\typing{\draw{gamma}\alpha{k,\lambda}}{\tyR m}}
\end{mathpar}
\caption{Subtyping and typing rules, with $m\in\{\G,\E\}$ and
$\alpha\in\{m,\M\}$. Rules for variables, unit and Boolean literals, functions,
recursion, application, let bindings, conditionals, pairs, projections, sum
injections and case analysis, and list constructors and case analysis are omitted
because they are standard.}
\label{fig:typing}
\Description{Subtyping permits G reals where E reals are expected. Multiplication requires a G left operand, division a G denominator, and comparison two G operands. Distribution rules share the same premises for sampling and mean evaluation.}
\end{figure}
